This report documents the design of the power stage of a three-phase
two-level voltage-source inverter for an automotive traction application:
\SI{800}{\volt} on the DC-link, \SI{30}{\Arms} per phase at any load power
factor, \SI{20}{\kilo\hertz} switching frequency, and forced-air or liquid
cooling with a maximum ambient of \SI{50}{\celsius}. The device is the
Wolfspeed \device{} silicon-carbide MOSFET~\cite{wolfspeed_e3m0016}, selected
against seven alternatives (other device classes, a paralleled pair, a
top-side-cooled package and a module) on criteria fixed before any result was
computed.

The work covers the design equations (operating point, DC-link capacitor,
modulation, dead time), a circuit-level PLECS model validated against
closed-form results, a SPICE study of the switching cell that fixes the gate
resistances, the loss and thermal budget, and the schematic and layout of the
power stage as a module of a larger board.

At the design point the converter delivers \SI{25.48}{\kilo\watt} of active
power and dissipates \SI{128.1}{\watt}, an efficiency of \SI{99.50}{\percent},
with a current ripple of \SI{0.84}{\percent} in RMS terms. The DC-link of
\SI{200}{\micro\farad} keeps the voltage ripple at \SI{0.17}{\percent} at the
rated point and at \SI{0.36}{\percent} in the worst operating condition,
against \SI{0.5}{\percent}. The junction settles at \SI{81.8}{\celsius}. The
design is sized by the stalled motor, a case this design adds to the brief:
with the full \SI{42.4}{\ampere} peak flowing as direct current through one
device, the junction reaches \SI{111.1}{\celsius} on an ideal cold plate
against a \SI{120}{\celsius} design criterion. On a real liquid-cooled rail,
computed in three dimensions, the criterion is met only after the insulating
interface is redesigned: with aluminium nitride and a phase-change material
the hottest junction reaches \SI{97.9}{\celsius}, and \SI{117.8}{\celsius}
with every tolerance stacked.

Four findings would each have produced a wrong report had they gone
unnoticed. The device first used was industrial rather than automotive and,
at locked rotor, twenty degrees beyond the criterion; the selection that
replaced it was decided by overvoltage at the die, not by temperature. The
first double-pulse bench measured a DC bus that was still charging, which hid
nearly \SI{70}{\volt} of overshoot; with the real bus the turn-off gate
resistance had to rise from \SI{5}{} to \SI{12.5}{\ohm}, and switching became
more than half of the loss budget. The vendor loss tables refer to the
body-diode free-wheeling case, while the library folder name suggests the
external-diode case: a factor of 2.0 on the turn-on energy. And the largest
thermal unknown is the insulating pad between the case and the cold plate,
more precisely its two interfaces, which on the real rail decide more of the
junction temperature than the coolant does.
